%=============================================================================
% Arknights Lone Trail LaTeX Template - Pure Vector TikZ Asset
% Asset: box_corner_mark.tex (Replacing 素材库/正文/definition_3.png, theorem_3.png, etc.)
% Original Dimensions: 69 x 69 bp
% Structure:
%   - Top-left & Bottom-right: hollow concentric rings (outer R=16.03, inner R=10.15)
%   - Bottom-left & Top-right: solid disks connected by smooth diagonal neck
% Dynamic Color Support:
%   - #1 defaults to 'definition'
%   - Accepts theorem, proposition, corollary, lemma, example, or custom colors
% Compatibility:
%   - Standalone compilation via xelatex (produces preview sheet)
%   - Direct %=============================================================================
% Arknights Lone Trail LaTeX Template - Pure Vector TikZ Asset
% Asset: box_corner_mark.tex (Replacing 素材库/正文/definition_3.png, theorem_3.png, etc.)
% Original Dimensions: 69 x 69 bp
% Structure:
%   - Top-left & Bottom-right: hollow concentric rings (outer R=16.03, inner R=10.15)
%   - Bottom-left & Top-right: solid disks connected by smooth diagonal neck
% Dynamic Color Support:
%   - #1 defaults to 'definition'
%   - Accepts theorem, proposition, corollary, lemma, example, or custom colors
% Compatibility:
%   - Standalone compilation via xelatex (produces preview sheet)
%   - Direct %=============================================================================
% Arknights Lone Trail LaTeX Template - Pure Vector TikZ Asset
% Asset: box_corner_mark.tex (Replacing 素材库/正文/definition_3.png, theorem_3.png, etc.)
% Original Dimensions: 69 x 69 bp
% Structure:
%   - Top-left & Bottom-right: hollow concentric rings (outer R=16.03, inner R=10.15)
%   - Bottom-left & Top-right: solid disks connected by smooth diagonal neck
% Dynamic Color Support:
%   - #1 defaults to 'definition'
%   - Accepts theorem, proposition, corollary, lemma, example, or custom colors
% Compatibility:
%   - Standalone compilation via xelatex (produces preview sheet)
%   - Direct %=============================================================================
% Arknights Lone Trail LaTeX Template - Pure Vector TikZ Asset
% Asset: box_corner_mark.tex (Replacing 素材库/正文/definition_3.png, theorem_3.png, etc.)
% Original Dimensions: 69 x 69 bp
% Structure:
%   - Top-left & Bottom-right: hollow concentric rings (outer R=16.03, inner R=10.15)
%   - Bottom-left & Top-right: solid disks connected by smooth diagonal neck
% Dynamic Color Support:
%   - #1 defaults to 'definition'
%   - Accepts theorem, proposition, corollary, lemma, example, or custom colors
% Compatibility:
%   - Standalone compilation via xelatex (produces preview sheet)
%   - Direct \input{box_corner_mark.tex} in documents (uses \boxcornercolor or definition)
%   - Direct macro call via \BoxCornerMark[<color>]
%=============================================================================
\ifcsname pgfpoint\endcsname
\else
  \documentclass[border=10pt]{standalone}
  \usepackage{tikz}
  \usepackage{xcolor}
  \def\imstandalone{1}
\fi

\providecolor{definition}{RGB}{245, 164, 0}
\providecolor{theorem}{RGB}{222, 4, 20}
\providecolor{proposition}{RGB}{177, 0, 129}
\providecolor{corollary}{RGB}{133, 194, 72}
\providecolor{lemma}{RGB}{233, 93, 157}
\providecolor{example}{RGB}{0, 153, 216}

\ifdefined\imstandalone
  \begin{document}
\fi

\providecommand{\BoxCornerMark}[1][definition]{%
  \begin{tikzpicture}[x=1bp, y=1bp, inner sep=0pt, outer sep=0pt]
    \useasboundingbox (0, 0) rectangle (69, 69);
    % Top-Left Ring: Center (16.03, 52.97), Outer R=16.03, Inner R=10.15
    \fill[#1, even odd rule] (16.03, 52.97) circle [radius=16.03] (16.03, 52.97) circle [radius=10.15];
    % Bottom-Right Ring: Center (52.97, 16.03), Outer R=16.03, Inner R=10.15
    \fill[#1, even odd rule] (52.97, 16.03) circle [radius=16.03] (52.97, 16.03) circle [radius=10.15];
    % Diagonal Dumbbell (BL to TR):
    \fill[#1] (16.03, 0.01)
      .. controls (7.19, 0.01) and (0.00, 7.19) .. (0.00, 16.03)
      .. controls (0.00, 24.87) and (7.19, 32.05) .. (16.03, 32.05)
      .. controls (16.92, 32.05) and (17.80, 31.98) .. (18.65, 31.84)
      .. controls (19.52, 31.70) and (20.40, 31.63) .. (21.27, 31.63)
      .. controls (25.50, 31.63) and (29.61, 33.30) .. (32.66, 36.35)
      .. controls (36.32, 40.01) and (38.00, 45.23) .. (37.16, 50.35)
      .. controls (37.01, 51.21) and (36.94, 52.09) .. (36.94, 52.97)
      .. controls (36.94, 61.81) and (44.13, 68.99) .. (52.97, 68.99)
      .. controls (61.81, 68.99) and (69.00, 61.81) .. (69.00, 52.97)
      .. controls (69.00, 44.13) and (61.81, 36.94) .. (52.97, 36.94)
      .. controls (52.08, 36.94) and (51.20, 37.02) .. (50.35, 37.16)
      .. controls (49.47, 37.30) and (48.59, 37.37) .. (47.72, 37.37)
      .. controls (43.49, 37.37) and (39.39, 35.70) .. (36.34, 32.65)
      .. controls (32.67, 28.99) and (30.99, 23.77) .. (31.84, 18.65)
      .. controls (31.98, 17.79) and (32.05, 16.91) .. (32.05, 16.03)
      .. controls (32.05, 7.19) and (24.87, 0.01) .. (16.03, 0.01) -- cycle;
  \end{tikzpicture}%
}

\ifcsname boxcornercolor\endcsname
\else
  \def\boxcornercolor{definition}
\fi

\makeatletter
\ifdefined\imstandalone
  \begin{tabular}{cccccc}
    \textbf{definition} & \textbf{theorem} & \textbf{proposition} & \textbf{corollary} & \textbf{lemma} & \textbf{example} \\
    \BoxCornerMark[definition] &
    \BoxCornerMark[theorem] &
    \BoxCornerMark[proposition] &
    \BoxCornerMark[corollary] &
    \BoxCornerMark[lemma] &
    \BoxCornerMark[example] \\[8pt]
    \resizebox{!}{8pt}{\BoxCornerMark[definition]} &
    \resizebox{!}{8pt}{\BoxCornerMark[theorem]} &
    \resizebox{!}{8pt}{\BoxCornerMark[proposition]} &
    \resizebox{!}{8pt}{\BoxCornerMark[corollary]} &
    \resizebox{!}{8pt}{\BoxCornerMark[lemma]} &
    \resizebox{!}{8pt}{\BoxCornerMark[example]} \\
    \small (8pt size) & \small (8pt size) & \small (8pt size) & \small (8pt size) & \small (8pt size) & \small (8pt size) \\
  \end{tabular}%
  \def\finishdoc{\end{document}}%
\else
  \ifx\@nodocument\relax
    \BoxCornerMark[\boxcornercolor]%
  \fi
  \let\finishdoc\relax%
\fi
\finishdoc
\makeatother

 in documents (uses \boxcornercolor or definition)
%   - Direct macro call via \BoxCornerMark[<color>]
%=============================================================================
\ifcsname pgfpoint\endcsname
\else
  \documentclass[border=10pt]{standalone}
  \usepackage{tikz}
  \usepackage{xcolor}
  \def\imstandalone{1}
\fi

\providecolor{definition}{RGB}{245, 164, 0}
\providecolor{theorem}{RGB}{222, 4, 20}
\providecolor{proposition}{RGB}{177, 0, 129}
\providecolor{corollary}{RGB}{133, 194, 72}
\providecolor{lemma}{RGB}{233, 93, 157}
\providecolor{example}{RGB}{0, 153, 216}

\ifdefined\imstandalone
  \begin{document}
\fi

\providecommand{\BoxCornerMark}[1][definition]{%
  \begin{tikzpicture}[x=1bp, y=1bp, inner sep=0pt, outer sep=0pt]
    \useasboundingbox (0, 0) rectangle (69, 69);
    % Top-Left Ring: Center (16.03, 52.97), Outer R=16.03, Inner R=10.15
    \fill[#1, even odd rule] (16.03, 52.97) circle [radius=16.03] (16.03, 52.97) circle [radius=10.15];
    % Bottom-Right Ring: Center (52.97, 16.03), Outer R=16.03, Inner R=10.15
    \fill[#1, even odd rule] (52.97, 16.03) circle [radius=16.03] (52.97, 16.03) circle [radius=10.15];
    % Diagonal Dumbbell (BL to TR):
    \fill[#1] (16.03, 0.01)
      .. controls (7.19, 0.01) and (0.00, 7.19) .. (0.00, 16.03)
      .. controls (0.00, 24.87) and (7.19, 32.05) .. (16.03, 32.05)
      .. controls (16.92, 32.05) and (17.80, 31.98) .. (18.65, 31.84)
      .. controls (19.52, 31.70) and (20.40, 31.63) .. (21.27, 31.63)
      .. controls (25.50, 31.63) and (29.61, 33.30) .. (32.66, 36.35)
      .. controls (36.32, 40.01) and (38.00, 45.23) .. (37.16, 50.35)
      .. controls (37.01, 51.21) and (36.94, 52.09) .. (36.94, 52.97)
      .. controls (36.94, 61.81) and (44.13, 68.99) .. (52.97, 68.99)
      .. controls (61.81, 68.99) and (69.00, 61.81) .. (69.00, 52.97)
      .. controls (69.00, 44.13) and (61.81, 36.94) .. (52.97, 36.94)
      .. controls (52.08, 36.94) and (51.20, 37.02) .. (50.35, 37.16)
      .. controls (49.47, 37.30) and (48.59, 37.37) .. (47.72, 37.37)
      .. controls (43.49, 37.37) and (39.39, 35.70) .. (36.34, 32.65)
      .. controls (32.67, 28.99) and (30.99, 23.77) .. (31.84, 18.65)
      .. controls (31.98, 17.79) and (32.05, 16.91) .. (32.05, 16.03)
      .. controls (32.05, 7.19) and (24.87, 0.01) .. (16.03, 0.01) -- cycle;
  \end{tikzpicture}%
}

\ifcsname boxcornercolor\endcsname
\else
  \def\boxcornercolor{definition}
\fi

\makeatletter
\ifdefined\imstandalone
  \begin{tabular}{cccccc}
    \textbf{definition} & \textbf{theorem} & \textbf{proposition} & \textbf{corollary} & \textbf{lemma} & \textbf{example} \\
    \BoxCornerMark[definition] &
    \BoxCornerMark[theorem] &
    \BoxCornerMark[proposition] &
    \BoxCornerMark[corollary] &
    \BoxCornerMark[lemma] &
    \BoxCornerMark[example] \\[8pt]
    \resizebox{!}{8pt}{\BoxCornerMark[definition]} &
    \resizebox{!}{8pt}{\BoxCornerMark[theorem]} &
    \resizebox{!}{8pt}{\BoxCornerMark[proposition]} &
    \resizebox{!}{8pt}{\BoxCornerMark[corollary]} &
    \resizebox{!}{8pt}{\BoxCornerMark[lemma]} &
    \resizebox{!}{8pt}{\BoxCornerMark[example]} \\
    \small (8pt size) & \small (8pt size) & \small (8pt size) & \small (8pt size) & \small (8pt size) & \small (8pt size) \\
  \end{tabular}%
  \def\finishdoc{\end{document}}%
\else
  \ifx\@nodocument\relax
    \BoxCornerMark[\boxcornercolor]%
  \fi
  \let\finishdoc\relax%
\fi
\finishdoc
\makeatother

 in documents (uses \boxcornercolor or definition)
%   - Direct macro call via \BoxCornerMark[<color>]
%=============================================================================
\ifcsname pgfpoint\endcsname
\else
  \documentclass[border=10pt]{standalone}
  \usepackage{tikz}
  \usepackage{xcolor}
  \def\imstandalone{1}
\fi

\providecolor{definition}{RGB}{245, 164, 0}
\providecolor{theorem}{RGB}{222, 4, 20}
\providecolor{proposition}{RGB}{177, 0, 129}
\providecolor{corollary}{RGB}{133, 194, 72}
\providecolor{lemma}{RGB}{233, 93, 157}
\providecolor{example}{RGB}{0, 153, 216}

\ifdefined\imstandalone
  \begin{document}
\fi

\providecommand{\BoxCornerMark}[1][definition]{%
  \begin{tikzpicture}[x=1bp, y=1bp, inner sep=0pt, outer sep=0pt]
    \useasboundingbox (0, 0) rectangle (69, 69);
    % Top-Left Ring: Center (16.03, 52.97), Outer R=16.03, Inner R=10.15
    \fill[#1, even odd rule] (16.03, 52.97) circle [radius=16.03] (16.03, 52.97) circle [radius=10.15];
    % Bottom-Right Ring: Center (52.97, 16.03), Outer R=16.03, Inner R=10.15
    \fill[#1, even odd rule] (52.97, 16.03) circle [radius=16.03] (52.97, 16.03) circle [radius=10.15];
    % Diagonal Dumbbell (BL to TR):
    \fill[#1] (16.03, 0.01)
      .. controls (7.19, 0.01) and (0.00, 7.19) .. (0.00, 16.03)
      .. controls (0.00, 24.87) and (7.19, 32.05) .. (16.03, 32.05)
      .. controls (16.92, 32.05) and (17.80, 31.98) .. (18.65, 31.84)
      .. controls (19.52, 31.70) and (20.40, 31.63) .. (21.27, 31.63)
      .. controls (25.50, 31.63) and (29.61, 33.30) .. (32.66, 36.35)
      .. controls (36.32, 40.01) and (38.00, 45.23) .. (37.16, 50.35)
      .. controls (37.01, 51.21) and (36.94, 52.09) .. (36.94, 52.97)
      .. controls (36.94, 61.81) and (44.13, 68.99) .. (52.97, 68.99)
      .. controls (61.81, 68.99) and (69.00, 61.81) .. (69.00, 52.97)
      .. controls (69.00, 44.13) and (61.81, 36.94) .. (52.97, 36.94)
      .. controls (52.08, 36.94) and (51.20, 37.02) .. (50.35, 37.16)
      .. controls (49.47, 37.30) and (48.59, 37.37) .. (47.72, 37.37)
      .. controls (43.49, 37.37) and (39.39, 35.70) .. (36.34, 32.65)
      .. controls (32.67, 28.99) and (30.99, 23.77) .. (31.84, 18.65)
      .. controls (31.98, 17.79) and (32.05, 16.91) .. (32.05, 16.03)
      .. controls (32.05, 7.19) and (24.87, 0.01) .. (16.03, 0.01) -- cycle;
  \end{tikzpicture}%
}

\ifcsname boxcornercolor\endcsname
\else
  \def\boxcornercolor{definition}
\fi

\makeatletter
\ifdefined\imstandalone
  \begin{tabular}{cccccc}
    \textbf{definition} & \textbf{theorem} & \textbf{proposition} & \textbf{corollary} & \textbf{lemma} & \textbf{example} \\
    \BoxCornerMark[definition] &
    \BoxCornerMark[theorem] &
    \BoxCornerMark[proposition] &
    \BoxCornerMark[corollary] &
    \BoxCornerMark[lemma] &
    \BoxCornerMark[example] \\[8pt]
    \resizebox{!}{8pt}{\BoxCornerMark[definition]} &
    \resizebox{!}{8pt}{\BoxCornerMark[theorem]} &
    \resizebox{!}{8pt}{\BoxCornerMark[proposition]} &
    \resizebox{!}{8pt}{\BoxCornerMark[corollary]} &
    \resizebox{!}{8pt}{\BoxCornerMark[lemma]} &
    \resizebox{!}{8pt}{\BoxCornerMark[example]} \\
    \small (8pt size) & \small (8pt size) & \small (8pt size) & \small (8pt size) & \small (8pt size) & \small (8pt size) \\
  \end{tabular}%
  \def\finishdoc{\end{document}}%
\else
  \ifx\@nodocument\relax
    \BoxCornerMark[\boxcornercolor]%
  \fi
  \let\finishdoc\relax%
\fi
\finishdoc
\makeatother

 in documents (uses \boxcornercolor or definition)
%   - Direct macro call via \BoxCornerMark[<color>]
%=============================================================================
\ifcsname pgfpoint\endcsname
\else
  \documentclass[border=10pt]{standalone}
  \usepackage{tikz}
  \usepackage{xcolor}
  \def\imstandalone{1}
\fi

\providecolor{definition}{RGB}{245, 164, 0}
\providecolor{theorem}{RGB}{222, 4, 20}
\providecolor{proposition}{RGB}{177, 0, 129}
\providecolor{corollary}{RGB}{133, 194, 72}
\providecolor{lemma}{RGB}{233, 93, 157}
\providecolor{example}{RGB}{0, 153, 216}

\ifdefined\imstandalone
  \begin{document}
\fi

\providecommand{\BoxCornerMark}[1][definition]{%
  \begin{tikzpicture}[x=1bp, y=1bp, inner sep=0pt, outer sep=0pt]
    \useasboundingbox (0, 0) rectangle (69, 69);
    % Top-Left Ring: Center (16.03, 52.97), Outer R=16.03, Inner R=10.15
    \fill[#1, even odd rule] (16.03, 52.97) circle [radius=16.03] (16.03, 52.97) circle [radius=10.15];
    % Bottom-Right Ring: Center (52.97, 16.03), Outer R=16.03, Inner R=10.15
    \fill[#1, even odd rule] (52.97, 16.03) circle [radius=16.03] (52.97, 16.03) circle [radius=10.15];
    % Diagonal Dumbbell (BL to TR):
    \fill[#1] (16.03, 0.01)
      .. controls (7.19, 0.01) and (0.00, 7.19) .. (0.00, 16.03)
      .. controls (0.00, 24.87) and (7.19, 32.05) .. (16.03, 32.05)
      .. controls (16.92, 32.05) and (17.80, 31.98) .. (18.65, 31.84)
      .. controls (19.52, 31.70) and (20.40, 31.63) .. (21.27, 31.63)
      .. controls (25.50, 31.63) and (29.61, 33.30) .. (32.66, 36.35)
      .. controls (36.32, 40.01) and (38.00, 45.23) .. (37.16, 50.35)
      .. controls (37.01, 51.21) and (36.94, 52.09) .. (36.94, 52.97)
      .. controls (36.94, 61.81) and (44.13, 68.99) .. (52.97, 68.99)
      .. controls (61.81, 68.99) and (69.00, 61.81) .. (69.00, 52.97)
      .. controls (69.00, 44.13) and (61.81, 36.94) .. (52.97, 36.94)
      .. controls (52.08, 36.94) and (51.20, 37.02) .. (50.35, 37.16)
      .. controls (49.47, 37.30) and (48.59, 37.37) .. (47.72, 37.37)
      .. controls (43.49, 37.37) and (39.39, 35.70) .. (36.34, 32.65)
      .. controls (32.67, 28.99) and (30.99, 23.77) .. (31.84, 18.65)
      .. controls (31.98, 17.79) and (32.05, 16.91) .. (32.05, 16.03)
      .. controls (32.05, 7.19) and (24.87, 0.01) .. (16.03, 0.01) -- cycle;
  \end{tikzpicture}%
}

\ifcsname boxcornercolor\endcsname
\else
  \def\boxcornercolor{definition}
\fi

\makeatletter
\ifdefined\imstandalone
  \begin{tabular}{cccccc}
    \textbf{definition} & \textbf{theorem} & \textbf{proposition} & \textbf{corollary} & \textbf{lemma} & \textbf{example} \\
    \BoxCornerMark[definition] &
    \BoxCornerMark[theorem] &
    \BoxCornerMark[proposition] &
    \BoxCornerMark[corollary] &
    \BoxCornerMark[lemma] &
    \BoxCornerMark[example] \\[8pt]
    \resizebox{!}{8pt}{\BoxCornerMark[definition]} &
    \resizebox{!}{8pt}{\BoxCornerMark[theorem]} &
    \resizebox{!}{8pt}{\BoxCornerMark[proposition]} &
    \resizebox{!}{8pt}{\BoxCornerMark[corollary]} &
    \resizebox{!}{8pt}{\BoxCornerMark[lemma]} &
    \resizebox{!}{8pt}{\BoxCornerMark[example]} \\
    \small (8pt size) & \small (8pt size) & \small (8pt size) & \small (8pt size) & \small (8pt size) & \small (8pt size) \\
  \end{tabular}%
  \def\finishdoc{\end{document}}%
\else
  \ifx\@nodocument\relax
    \BoxCornerMark[\boxcornercolor]%
  \fi
  \let\finishdoc\relax%
\fi
\finishdoc
\makeatother

